\Needspace{8\baselineskip}
\Article{197}
If we have one transformation of a form $(a,b,c)\mathinner{\ldotp\ldotp\ldotp}\varphi$ into an equivalent form $\Phi$, we can derive from it \textit{all} similar transformations of $\varphi$ into $\Phi$, provided we can find all the solutions of the indeterminate equation $tt-Duu=mm$,
\SourceLeaf{192}
where \(D\) denotes the determinant of the forms \(\Phi,\varphi\), and \(m\) the greatest common divisor of \(a,2b,c\) (Article~162). Having already solved this problem above for a negative value of \(D\), we will now tackle it for a positive value. Clearly, any value of \(t\) satisfying the equation still satisfies it when its sign is reversed, and the same is true of any value of \(u\). So it is enough to find all \textit{positive} values of \(t,u\), and each solution in positive values will stand for four solutions. We will do this by showing \textit{first} how to find the \textit{smallest} values of \(t,u\) (apart from the obvious ones \(t=m,\ u=0\)), and \textit{then} how to derive all the others from them.

\Needspace{8\baselineskip}
\Article{198}
\textsc{Problem.} \textit{Find the smallest numbers \(t,u\) satisfying the indeterminate equation \(tt-Duu=mm\), given a form \((M,N,P)\) whose determinant is \(D\) and for which \(M,2N,P\) have greatest common divisor \(n\).}

\textit{Solution.} Choose any reduced form \((a,b,-a')\ldots f\) with determinant \(D\), in which \(a,2b,a'\) have greatest common divisor \(m\). Such a form clearly exists, since we can find a reduced form equivalent to \((M,N,P)\), which by Article~161 has this property. But for our present purpose we can use any reduced form satisfying this condition. Work out the period of \(f\), which we will suppose consists of \(n\) forms. Retaining all the notation used in Article~188, \(f^n\) is \((+a^n,b^n,-a^{n+1})\), since \(n\) is even; and \(f\) passes into this form by the proper substitution \(\alpha^n,\beta^n,\gamma^n,\delta^n\). But \(f\) and \(f^n\) are identical, so \(f\) also passes into \(f^n\) by the proper substitution \(1,0,0,1\). From these two similar transformations of \(f\) into \(f^n\), Article~162 gives a solution of \(tt-Duu=mm\) \textit{in integers}, namely \(t=\frac12(\alpha^n+\delta^n)m\) (equation~18, Article~162), \(u=\frac{\gamma^n m}{a}\) (equation~19)\HistoricalNote{*)}{\begin{tabular}[t]{@{}l@{}}\(\begin{array}{@{}rll@{}}
\text{In Article 162 we had}&\alpha,\beta,\gamma,\delta;&\alpha',\beta',\gamma',\delta';\\
\text{here we have}&1,0,0,1;&\alpha^n,\beta^n,\gamma^n,\delta^n;\end{array}\)
\\[0.6ex]\(\begin{array}{@{}lll@{}}A,B,C;&a,b,c;&e,\\a,b,-a';&a,b,-a';&1.\end{array}\)\end{tabular}}. Denote these values by \(T,U\), taking them positive if they are not already so. These \(T,U\) will be the smallest values of \(t,u\), apart from \(t=m,\ u=0\) (from which they must differ, since \(\gamma^n\) clearly cannot be \(=0\)).

For suppose there are still smaller values of \(t,u\), say \(\mathfrak{t},\mathfrak{u}\), which are positive, with \(\mathfrak{u}\) not \(=0\). Then by Article~162 the form \(f\) is transformed by the proper substitution \(\frac1m(\mathfrak{t}-b\mathfrak{u}),\ \frac1m a'\mathfrak{u},\ \frac1m a\mathfrak{u},\ \frac1m(\mathfrak{t}+b\mathfrak{u})\) into a form identical to
\SourceLeaf{193}
itself. By Article~193, II, \textit{either} \(\frac1m(\mathfrak{t}-b\mathfrak{u})\) \textit{or} \(-\frac1m(\mathfrak{t}-b\mathfrak{u})\) must equal one of the numbers \(\alpha'',\alpha''',\alpha''''\) etc., say \(=\alpha^\mu\) (for since \(\mathfrak{tt}=D\mathfrak{uu}+mm=bb\mathfrak{uu}+aa'\mathfrak{uu}+mm\), we have \(\mathfrak{tt}>bb\mathfrak{uu}\), so \(\mathfrak{t}-b\mathfrak{u}\) is positive; hence the fraction \(\frac{\mathfrak{t}-b\mathfrak{u}}{a\mathfrak{u}}\), corresponding to \(\frac{\mathfrak{A}}{\mathfrak{C}}\) in Article~193, has the same sign as \(a\) or \(a'\)). In the first case \(\frac1m a'\mathfrak{u},\ \frac1m a\mathfrak{u},\ \frac1m(\mathfrak{t}+b\mathfrak{u})\), and in the second the same quantities with their signs reversed, are respectively \(=\beta^\mu,\gamma^\mu,\delta^\mu\). But since \(\mathfrak{u}<U\), \textit{i.e.}, \(\mathfrak{u}<\frac{\gamma^n m}{a}\) and \(>0\), we have \(\gamma^\mu<\gamma^n\) and \(>0\). Since the sequence \(\gamma,\gamma',\gamma''\) etc. increases continually, \(\mu\) must lie strictly between \(0\) and \(n\). But the corresponding form \(f^\mu\) is identical to \(f\), \textit{a contradiction}, since all the forms \(f,f',f''\) etc. as far as \(f^{n-1}\) are assumed distinct. We conclude that the smallest values of \(t,u\) (apart from \(m,0\)) are \(T,U\).

\textit{Example.} If \(D=79,\ m=1\), we can use the form \((3,8,-5)\), for which \(n=6\) and \(\alpha^n=-8,\ \gamma^n=-27,\ \delta^n=-152\) (Article~188). This gives \(T=80,\ U=9\), the smallest values of \(t,u\) satisfying \(tt-79uu=1\).

\Needspace{8\baselineskip}
\Article{199}
For practical use we can obtain still more convenient formulas. We have \(2b\gamma^n=-a(\alpha^n-\delta^n)\), which follows easily from Article~162 by multiplying equation~\([19]\) by \(2b\) and \([20]\) by \(a\), and changing the notation used there to our present notation. This gives \(\alpha^n+\delta^n=2\delta^n-\frac{2b}{a}\gamma^n\), so
\[
\pm T=m\left(\delta^n-\frac{b}{a}\gamma^n\right),\quad \pm U=\frac{\gamma^n m}{a}
\]
By a similar method we obtain
\[
\pm T=m\left(\alpha^n+\frac{b}{a'}\beta^n\right),\quad \pm U=\frac{\beta^n m}{a'}
\]
Both sets of formulas are very convenient, since \(\gamma^n=\delta^{n-1}\), \(\alpha^n=\beta^{n-1}\). If we use the latter, it is enough to calculate just the sequence \(\beta',\beta'',\beta'''\ldots\beta^n\); if we prefer the former, just the sequence \(\delta',\delta'',\delta'''\) etc. Further, Article~189, 3 easily shows that, since \(n\) is necessarily even, \(\alpha^n\) and \(\frac{b}{a'}\beta^n\) have the same sign, as do \(\delta^n\) and \(\frac{b}{a}\gamma^n\). So in the first formula for \(T\) we must take the absolute difference, and in the second the absolute sum; we need not consider the signs at all. Returning to the notation used in Article~189, 4, the first formula gives
\[
\begin{aligned}T&=m[k',k'',k'''\ldotp\,\ldotp\,\ldotp\,\ldotp\, k^n]-\frac{mb}{a}[k',k'',k'''\ldotp\,\ldotp\,\ldotp\,\ldotp\, k^{n-1}],\\U&=\frac{m}{a}[k',k'',k'''\ldotp\,\ldotp\,\ldotp\,\ldotp\, k^{n-1}]\end{aligned}
\]
\SourceLeaf{194}
and the second
\[
\begin{aligned}T&=m[k'',k'''\ldotp\,\ldotp\,\ldotp\,\ldotp\, k^{n-1}]+\frac{mb}{a'}[k'',k'''\ldotp\,\ldotp\,\ldotp\,\ldotp\, k^n],\\U&=\frac{m}{a'}[k'',k'''\ldotp\,\ldotp\,\ldotp\,\ldotp\, k^n]\end{aligned}
\]
where the value of \(T\) can also be written \(m\left[k'',k'''\ldotp\,\ldotp\,\ldotp\,\ldotp\, k^n,\frac{b}{a'}\right]\).

\textit{Example.} For \(D=61,\ m=2\) we can use the form \((2,7,-6)\), for which we find \(n=6\), and \(k',k'',k''',k'''',k''''',k''''''\) respectively \(=2,2,7,2,2,7\). This gives
\[
\begin{aligned}T&=2[2,2,7,2,2,7]-7[2,2,7,2,2]\\&=2888-1365=1523\end{aligned}
\]
from the first formula; the second gives the same result:
\[
T=2[2,7,2,2]+\tfrac73[2,7,2,2.7]
\]
\[
U\text{ is }=[2,2,7,2,2]=\tfrac13[2,7,2,2,7]=195.
\]
There are other ways to shorten the calculation, but to keep this account brief we cannot discuss them more fully here.

\Needspace{8\baselineskip}
\Article{200}
To obtain \textit{all} values of \(t,u\) from the smallest ones, we write the equation \(TT-DUU=mm\) as
\[
\left(\frac{T}{m}+\frac{U}{m}\surd D\right)\left(\frac{T}{m}-\frac{U}{m}\surd D\right)=1
\]
which also gives
\[
\left(\frac{T}{m}+\frac{U}{m}\surd D\right)^e\left(\frac{T}{m}-\frac{U}{m}\surd D\right)^e=1\quad\mathinner{\ldotp\ldotp\ldotp\ldotp\ldotp\ldotp\ldotp}\tag*{[1]}
\]
where \(e\) denotes any number. For brevity, we will generally denote the values of the quantities
\begingroup\postdisplaypenalty=10000
\[
\begin{gathered}\frac{m}{2}\left(\frac{T}{m}+\frac{U}{m}\surd D\right)^e+\frac{m}{2}\left(\frac{T}{m}-\frac{U}{m}\surd D\right)^e,\\
\frac{m}{2\surd D}\left(\frac{T}{m}+\frac{U}{m}\surd D\right)^e-\frac{m}{2\surd D}\left(\frac{T}{m}-\frac{U}{m}\surd D\right)^e\text{\textsuperscript{*)}}\end{gathered}
\]
\HistoricalDisplayNote{*)}{Only in these four expressions and in equation~\([1]\) does \(e\) denote the \textit{exponent} of a power; elsewhere superscript letters always denote an \textit{index}.}\endgroup
by \(t^e,u^e\) respectively: \textit{i.e.}, their values for \(e=0\) by \(t^0,u^0\) (which are \(m,0\)); for \(e=1\) by \(t',u'\) (which are \(T,U\)); for \(e=2\) by \(t'',u''\); for \(t=3\) by \(t''',u'''\) etc. --- We will prove that if \(e\) takes all nonnegative integers, \textit{i.e.}, \(0\) and all positive integers from \(1\) to \(\infty\), these expressions give all positive values of \(t,u\): namely, I) all values of these expressions
\SourceLeaf{195}
are indeed values of \(t,u\); II) all these values are integers; III) there are no positive values of \(t,u\) not contained in these formulas.

I. Substituting the values of \(t^e,u^e\), we readily confirm, using equation~\([1]\), that
\[
(t^e+u^e\surd D)(t^e-u^e\surd D)=mm,\quad \textit{i.e.}\quad t^e t^e-Du^e u^e=mm
\]

II. In the same way we easily confirm that, in general,
\[
t^{e+1}+t^{e-1}=\frac{2T}{m}t^e,\quad u^{e+1}+u^{e-1}=\frac{2T}{m}u^e
\]
The two sequences \(t^0,t',t'',t'''\) etc., \(u^0,u',u'',u'''\) etc. therefore satisfy recurrence relations, both with coefficients \(\frac{2T}{m},-1\), namely
\[
t''=\frac{2T}{m}t'-t^0,\quad t'''=\frac{2T}{m}t''-t'\text{ etc.,}\quad u''=\frac{2T}{m}u'\text{ etc.}
\]
Now by hypothesis there is a form \((M,N,P)\) with determinant \(D\), in which \(M,2N,P\) are divisible by \(m\). So
\[
TT=(NN-MP)UU+mm
\]
and \(4TT\) is clearly divisible by \(mm\). Thus \(\frac{2T}{m}\) is an integer, and a positive one. But we have \(t^0=m,\ t'=T,\ u^0=0,\ u'=U\), all integers, so all of \(t'',t'''\) etc., \(u'',u'''\) etc. are integers too. Further, since \(TT>mm\), all of \(t^0,t',t'',t'''\) etc. are clearly positive and increase continually without bound, as do all of \(u^0,u',u'',u'''\) etc.

III. Suppose there are other positive values of \(t,u\) not contained in the sequences \(t^0,t',t''\) etc., \(u^0,u',u''\) etc., say \(\mathfrak{T},\mathfrak{U}\). Since the sequence \(u^0,u'\) etc. increases from \(0\) without bound, \(\mathfrak{U}\) must clearly lie between two consecutive terms, \(u^n\) and \(u^{n+1}\), so that \(\mathfrak{U}>u^n\) and \(\mathfrak{U}<u^{n+1}\). To show that this assumption is absurd, we observe the following.

\(1^\circ\) The equation \(tt-Duu=mm\) is also satisfied by taking
\[
t=\frac1m(\mathfrak{T}t^n-D\mathfrak{U}u^n),\quad u=\frac1m(\mathfrak{U}t^n-\mathfrak{T}u^n).
\]
\SourceLeaf{196}
This is readily confirmed by substitution. We show as follows that these values, which we will briefly denote by \(=\tau,\upsilon\), are always \textit{integers}. If \((M,N,P)\) is a form with determinant \(D\), and \(m\) is a common divisor of \(M,2N,P\), then both \(\mathfrak{T}+N\mathfrak{U}\) and \(t^n+Nu^n\) are divisible by \(m\), and so is \(\mathfrak{U}(t^n+Nu^n)-u^n(\mathfrak{T}+N\mathfrak{U})\), or \(\mathfrak{U}t^n-\mathfrak{T}u^n\). Thus \(\upsilon\) is an integer, and so is \(\tau\), since \(\tau\tau=D\upsilon\upsilon+mm\).

\(2^\circ\) Clearly \(\upsilon\) cannot be \(=0\), for this would give
\[
\mathfrak{UU}t^n t^n=\mathfrak{TT}u^n u^n
\]
or
\[
\mathfrak{UU}(Du^n u^n+mm)=u^n u^n(D\mathfrak{UU}+mm)
\]
or \(\mathfrak{UU}=u^n u^n\), contrary to the hypothesis \(\mathfrak{U}>u^n\). So, since apart from \(0\) the smallest value of \(u\) is \(U\), \(\upsilon\) is certainly no smaller than \(U\).

\(3^\circ\) From the values of \(t^n,t^{n+1},u^n,u^{n+1}\) we easily confirm that
\[
mU=u^{n+1}t^n-t^{n+1}u^n.
\]
So \(\mathfrak{U}t^n-\mathfrak{T}u^n\) is certainly no smaller than \(u^{n+1}t^n-t^{n+1}u^n\).

\(4^\circ\) Now the equation \(\mathfrak{TT}-D\mathfrak{UU}=mm\) gives
\[
\frac{\mathfrak{T}}{\mathfrak{U}}=\surd\left(D+\frac{mm}{\mathfrak{UU}}\right)
\]
and similarly
\[
\frac{t^{n+1}}{u^{n+1}}=\surd\left(D+\frac{mm}{u^{n+1}u^{n+1}}\right)
\]
from which we easily deduce \(\frac{\mathfrak{T}}{\mathfrak{U}}>\frac{t^{n+1}}{u^{n+1}}\). This and the conclusion in \(3^\circ\) give
\[
\begin{aligned}&(\mathfrak{U}t^n-\mathfrak{T}u^n)\left(t^n+u^n\frac{\mathfrak{T}}{\mathfrak{U}}\right)\\&\qquad>(u^{n+1}t^n-t^{n+1}u^n)\left(t^n+u^n\frac{t^{n+1}}{u^{n+1}}\right)\end{aligned}
\]
or, expanding and substituting for \(\mathfrak{TT},\ t^n t^n,\ t^{n+1}t^{n+1}\) their values \(D\mathfrak{UU}+mm,\allowbreak\ Du^n u^n+mm,\allowbreak\ Du^{n+1}u^{n+1}+mm\),
\[
\frac1{\mathfrak{U}}(\mathfrak{UU}-u^n u^n)>\frac1{u^{n+1}}(u^{n+1}u^{n+1}-u^n u^n)
\]
Since both quantities are clearly positive, rearrangement gives \(\mathfrak{U}+\frac{u^n u^n}{u^{n+1}}>u^{n+1}+\frac{u^n u^n}{\mathfrak{U}}\), \textit{a contradiction}, because the first part of the first quantity is \textit{smaller}
\SourceLeaf{197}
than the first part of the second quantity, and its second part is also smaller than the second part of the other. So the assumption cannot hold, and the sequences \(t^0,t',t''\) etc., \(u^0,u',u''\) etc. give all positive values of \(t,u\).

\textit{Example.} For \(D=61,\ m=2\) we found the smallest positive values of \(t,u\) to be \(1523,195\). So \textit{all} positive values are given by the formulas
\[
\begin{aligned}
t&=\left(\frac{1523}{2}+\frac{195}{2}\surd61\right)^e+\left(\frac{1523}{2}-\frac{195}{2}\surd61\right)^e\\
u&=\frac1{\surd61}\left(\left(\frac{1523}{2}+\frac{195}{2}\surd61\right)^e-\left(\frac{1523}{2}-\frac{195}{2}\surd61\right)^e\right)
\end{aligned}
\]
We find
\[
\begin{aligned}
t^0&=2,\quad t'=1523,\\
t''&=1523t'-t^0=2319527,\\
t'''&=1523t''-t'=3532618098\text{ etc.}\\
u^0&=0,\quad u'=195,\\
u''&=1523u'-u^0=296985,\\
u'''&=1523u''-u'=452307960\text{ etc.}
\end{aligned}
\]
